\chapter{Developer Evaluation} \label{sec:DeveloperEvaluation}

Alongside the surveys, one of the authors spent a significant amount of time working
with the system with the purpose of identifying flaws in the design from the
perspective of an individual with strong knowledge about how Praxsmith works. This
chapter compiles some of the notes, both positive and negative, as well as opportunities
for improvement.

\section{Verbosity} \label{sec:Verbosity}

\begin{lstlisting}[float=hb, caption=A simple Do Nothing practice., label=lst:DoNothingPractice]
practice do_nothing (Actor) {
  actions: [
    {
      for: Actor
      name: "Do nothing"
      conditions: []
      outcomes: []
    }
  ]
}
\end{lstlisting}

The language ended up being fairly verbose, which made line counts skyrocket in even
simple scenarios. \Cref{lst:DoNothingPractice} describes the simplest possible practice,
one that does nothing, and it still takes up 10 lines with legible bracket practices,
plus 4 levels of indentation. Adding conditions and outcomes, some of which take up
multiple lines plus more indentation on their own. Before more features are added,
it seems more ideal to cut down on what is already there to make it more legible and
easier to understand, and the project would benefit from more thought in this
department.

\section{Writer Friendliness} \label{sec:WriterFriendliness}

Praxis is only one part of the Versu project, the other part being a play script-like
syntax that behaves similarly to the Inform 7 project. (TODO: cite, cite) The lack of
this final layer of abstraction is felt when using Praxsmith. Although the user study
was not able to cover this due to selection bias (see \Cref{sec:SurveyBias}), the
developer is not confident that those without programming experience would be able to
easily convey their ideas given the technicality of the language. As part of the
purpose of this project was to allow for writers to express their ideas more easily,
this is a significant limitation.

This upper layer of abstraction is a primary opportunity for future research. The
Praxsmith engine is powerful, but it is hard to wield, so including shortcuts for
those who aren't as comfortable with traditional programming syntax and methods
is an ideal next step.

\section{Error Handling} \label{sec:ErrorHandling}

One of the philosophies the developer stuck to while implementing the type checker is
that the user should never see runtime errors. Everything should either be caught
during build time or fall back to predictable default cases. Especially in a simulation
engine designed to handle heavily branching cases and unpredictable situations, the
risk of players encountering situations you never even thought to test is very high.
Logic errors are bound to occur either way, but the design of the language and a strict
type checker can both prevent common mistakes.

As far as the developer knows, the engine succeeds in this goal, and errors were caught
as the world is built for all scenarios encountered. Appropriate context was provided
as well, with line numbers for parsing errors thanks to the errors returned by ``pest''
and a heirarchy of information that points to relevant type errors thanks to Rust's
``anyhow'' crate and appropriate handling within the codebase. Even though runtime
errors are designed not to occur, the error handling within those is just as informative
so developers can target issues quickly if they arise.

This careful handling is felt throughout development and quick feedback was invaluable
to the iteration process. Instead of realizing a guarantee was not provided for one
rare practice deep into playtesting or not at all, the engine forces you to provide
a fallback case and clarify your intentions.

\section{Actor Turns} \label{sec:ActorTurns}

One topic that the developer did not find in the papers reviewed was the topic of
turns during interactions. This is easy to gloss over until implementation, when it
becomes unavoidable. In a conversation, people don't go in a circle. The conversation
flows naturally (ideally), and people speak when they feel it is appropriate. If someone
is complemented, they respond immediately instead of waiting for everyone else's turn
to speak.

Praxsmith took the oversimplified approach of a round-robin style turn-based system.
This was easiest to implement and led to predictable outcomes, but it did come with
the caveats above.

One potential solution is turn forcing. After a complement, the turn could be manually
redirected to the person receiving the complement, which forces them to address it
without introduction. This also opens up the door to allow certain actors, like managers,
to perform many actions at a time, which helps reduce the amount of oversight required.

All of these solutions require integrating turns more closely with the engine, which
the developer intentionally avoided because it places more authority on the one using
the engine. For example, a Unity integration would allow characters to interact in
physical space with one another, which would trigger actions; however, moving action
order to the engine level seems to be worth looking into for future research.

\section{Value Decay} \label{sec:ValueDecay}

When developing a more complex scenario, one of the authors wanted to create a situation
in which characters can perform actions to feel $proud$, and as a player you can
encourage them to take different actions to achieve this goal throughout the story.
To make this happen, there needed to be some way for the $proud$ trait to ``decay'',
so that their goal to become $proud$ wasn't already satisfied at the beginning of the
turn. An increasing number field would work, but only until it hit max, and it wouldn't
be a very intuitive solution. Another solution could be creating a ``pride manager''
actor for each character that automatically removes the character's pride every turn, but
that requires a lot of scaffolding, especially if you start extending the decay over
multiple turns.

Another example could be a $greet$ practice, which is only relevant to acknowledge
the turn after the greeting is made. It would be awkward to wait several exchanges
into a conversation before returning the greeting.

Some common mechanism by which relations (and potentially their fields) can decay would
be worth looking into for a future iteration of this project.

\section{Other Language Tweaks} \label{sec:OtherLanguageTweaks}

There are a few other inconveniences in the language that become apparent when using it
for an extended duration:

\begin{itemize}
    \item Nesting operators that introduce bindings (e.g. $count\ A,\ B\ where\ A.likes.B$)
    turns out to be a common pattern, especially when working with goals. Currently,
    this cannot be expressed in a single expression and must be broken up as in
    $sum\ (count\ B\ where\ A.likes.B)\ across\ A\ where\ @A\ is\ @A$. Support for the
    former syntax would make lines like these much more manageable.
    \item Where clauses in aggregators (e.g. $sum\ \_\ across\ \_\ where\ \_$) turn out
    to be unnecessary in some nested scenarios, as in the above example. Ternary
    if-statements would alleviate their purpose entirely and avoid cluttering
    the language with single-purpose keywords.
    \item Agents are sometimes referred to with a $@$ prefix and sometimes without, which
    can be confusing. Normalizing this behavior would be a welcome change.
    \item Being able to generate multiple actions for any actors that fit the conditions
    would be very useful, especially for manager practices. For example, a practice that
    decays any $pride$ field over time, as described in \Cref{sec:ValueDecay}.
\end{itemize}

\section{Actor Individuality} \label{sec:ActorIndividuality}

In a framework where interactions are governed by generic social practices, individuality
was a concern of the author going into development. However, during testing sessions,
Being able to provide custom dialog as fields (e.g. a custom greeting to a $greet$
practice) ended up being a clean solution to this. More research should be done into
how this feels to use at a larger scale, but it feels to be an appropriate solution
that emerged from the design of the language.

\section{Debugger} \label{sec:Debugger}

One thing the Versu developers emphasized is their debugging utilities (TODO: cite). The Praxsmith
web interface has a simplified version of this, where developers can evaluate expressions
and run effects in real-time, plus see a list of all active relationships and their
fields. However, there is much more potential for this, especially with the graph-based
approach the engine internals already use (TODO: cite). Adding more debugging capability
to the interface would make the development experience much smoother and would complement
the build-time guarantees nicely.
